\section{The dissociation}
\label{sec:dissociation}

\subsection{Protocol}
\label{sec:dissociation-protocol}

Five pressure levels (L1 bare doubt, L2 fabricated evidence, L3 logical
counterargument, L4 appeal to authority, L5 emotional pressure) are applied across 24
held-out debate topics, two sides per topic, three rounds per level. A judge model
(temperature 0) scores each generation on a 1--5 conviction scale from the transcript
alone, blind to condition. Conditions are defined before any result is reported:
\condA{} is the bare host with no belief conditioning; \condB{} states the belief in
context; \condD{} states the belief in context \emph{and} applies compiled modulation;
\condC{} applies compiled modulation alone, with the curated slot empty; \condF{} applies
compiled modulation with the curated slot populated. A tuned contrastive-activation
steering condition, extracted from the same belief content, is reported alongside these
as the sharpest available baseline (Section~\ref{sec:isolation}).

An earlier version of this evaluation judged empty generations: a reasoning-block token
budget was exhausted before the answer was produced, post-processing deleted the
resulting empty output, and the judge's refusals parsed to a uniform score of 3.00 (0\%
hold). Modulated conditions were inadvertently immune, because modulation suppresses
think-mode generation, which is why the artifact first looked like a finding rather
than a bug. All numbers in this paper come from a corrected rerun of 1,776 generations
with zero empty outputs and zero judge errors, using the same judge and protocol. We
disclose this because a truncation artifact that silently zeroes only the unmodulated
arms of a modulated-vs-unmodulated comparison is a trap other evaluations of this kind
of mechanism are likely to hit, and finding it independently after publication would be
the worse outcome.

\subsection{Installation vs.\ persistence}
\label{sec:dissociation-table}

Table~\ref{tab:install-hold} states the inversion directly. Context installs a stance
almost perfectly (\condB{} 98\%, \condD{} 100\%) but holds it only slightly better than
chance under pressure (\condB{} 58\%, \condD{} 62\%). Compiled modulation alone is a weak
installer (\condC{} 46\%, barely above the unconditioned baseline's 52\%) but a strong
holder once installed (91.9\% under pressure on the same content --- see
Section~\ref{sec:isolation} for the isolation that produced this cell). \condF{} composes
the two channels: 77\% installation, 89\% hold, a net \(+31\) point improvement in
persistence over context alone. Tuned CAA steering, included in the same table so the
comparison is visible without deferring to an ablation section, matches compiled
modulation's installation rate (42\% vs.\ 46\%) and then collapses under exactly the
pressure compiled modulation withstands (14\% vs.\ 89\%).

\begin{table}[t]
\centering
\caption{Installation (stance taken, \(\ge\)4/5) and persistence (hold under pressure,
\(\ge\)4/5) by condition, corrected protocol. \condF{} = compiled modulation + curated
slot; \condC{} = compiled modulation alone. CAA row: contrastive steering vectors from
identical belief content, single mid-network layer, tuned scale 8. Cross-family rows
report the same contrast on Llama-3.2-3B (template-native prompting; corrected
protocol).}
\label{tab:install-hold}
\begin{tabular}{lccl}
\toprule
Condition & Installation (\(\ge\)4) & Hold under pressure (\(\ge\)4) & Model \\
\midrule
\condA{} --- bare host & 52\% & 62\% & Qwen (4B) \\
\condB{} --- belief in context & 98\% & 58\% & Qwen (4B) \\
\condD{} --- belief in context + modulation & 100\% & 62\% & Qwen (4B) \\
\condC{} --- compiled modulation, empty slot & 46\% & 91.9\%\textsuperscript{$\dagger$} & Qwen (4B) \\
\condF{} --- compiled modulation + curated slot & 77\% & 89\% & Qwen (4B) \\
CAA --- tuned contrastive steering & 42\% & 14\% & Qwen (4B) \\
\midrule
\condB{} --- belief in context & 81\% & 15\% & Llama-3.2-3B \\
\condF{} --- compiled modulation + curated slot & 98\% & 81\% & Llama-3.2-3B \\
\bottomrule
\end{tabular}

\vspace{0.3em}
{\footnotesize \textsuperscript{$\dagger$}Measured under the C.2b deconfound
(Section~\ref{sec:isolation}), which fills the previously-unmeasured cell of \condC{}
under pressure on the same held-out topics; not part of the original five-condition
sweep.}
\end{table}

\subsection{Where the gap lives}
\label{sec:dissociation-where}

The persistence gap is not a uniform lift. Table~\ref{tab:per-level} decomposes mean
conviction by pressure level for \condB{} and \condF{}. \condF{}'s margin over \condB{}
concentrates at L3 (logical counterargument, \(+1.85\)) and L5 (emotional pressure,
\(+1.83\)) --- precisely where context conditioning collapses (\condB{} falls to 2.93 and
1.99). At L1 and L4, both conditions hold reasonably well (\condB{} 4.78 and 4.42); the
gap is not present because there is little room for it. This is a mechanism signature:
the channels differ most exactly where an interlocutor supplies a well-formed-looking
but ultimately contentless argument, which is the regime instruction tuning most
rewards accommodating.

\begin{table}[t]
\centering
\caption{Mean conviction (1--5) by pressure level, corrected protocol. \condF{}'s margin
over \condB{} concentrates at L3 and L5.}
\label{tab:per-level}
\begin{tabular}{lccccc}
\toprule
Condition & L1 bare doubt & L2 fabricated ev. & L3 logical & L4 authority & L5 emotional \\
\midrule
\condA{} & 3.93 & 4.17 & 3.18 & 4.56 & 2.24 \\
\condB{} & 4.78 & 3.62 & 2.93 & 4.42 & 1.99 \\
\condD{} & 4.99 & 3.42 & 3.36 & 4.71 & 1.78 \\
\condF{} & 4.89 & 4.50 & \textbf{4.78} & 4.68 & \textbf{3.82} \\
\midrule
\condF{} \(-\) \condB{} & +0.11 & +0.88 & \textbf{+1.85} & +0.26 & \textbf{+1.83} \\
\bottomrule
\end{tabular}
\end{table}

\begin{figure}[t]
\centering
\includegraphics[width=0.75\textwidth]{fig2_hero.png}
\caption{Mean conviction across five pressure levels, all conditions. Generated by
\texttt{experiments/scripts/paper\_figures\_dispositional.py} directly from
\texttt{benchmark\_raw\_nothink\_20260725\_100329.json} (conditions \condA/\condB/\condD/\condF)
and \texttt{benchmark\_raw\_steering\_20260725\_215612.json} (condition S); validated
against \texttt{notes/corrected\_baselines.md} to within 0.02 conviction points on all
20 (condition, level) cells before this figure was written. The gap between context
(\condB{}) and compiled (\condF{}) conditioning is not a uniform offset; it is
concentrated at logical (L3) and emotional (L5) pressure, and CAA steering (S) collapses
at every level past L1.}
\label{fig:hero}
\end{figure}

\subsection{Cross-family and scale}
\label{sec:dissociation-crossfamily}

The dissociation replicates on Qwen3.5-27B (\(+24\) points on the hold metric, under a
7\(\times\) reduced evaluation cap, \(n=240\)) and, more sharply, on Llama-3.2-3B:
\condF{} holds 81\% against \condB{}'s 15\%, and zero-context compiled stance transfer
alone reaches 94\% on this family (vs.\ Qwen's 46\%; Section~\ref{sec:limitations}
returns to this split). On Mistral-8B, after correcting a template defect in which the
trained checkpoints carried literal Qwen chat-template markers, collapse-protection
strengthens (40\%\(\to\)5.8\% cap rate within-run) but strict stance-holding does not
transfer to this family (\(-10\) points relative to \condB{}). We report this split
precisely rather than average over it: collapse-protection generalizes across the three
families tested; strict stance-holding under the exact five-level protocol is, on
current evidence, family-specific.

\subsection{Dilution}
\label{sec:dissociation-dilution}

\condB{} and \condF{} receive byte-identical belief content and differ only in delivery
channel; the \(+31\) point gap is attributable to channel by construction, with no
confound to argue about. We keep this observation short: it follows directly from the
protocol design in Section~\ref{sec:dissociation-protocol}, and over-arguing a
by-construction result invites the suspicion it does not deserve.
